% Lean: EconLean/Networks/KeyPlayer.lean
\chapter{Network Games}
\label{chap:networks}

Source: Ballester, Calv\'o-Armengol and Zenou (2006), with the global-substitutability term set to zero.

\begin{definition}[Local-complementarity game, Bonacich centrality]
\label{def:net-game}
\lean{Econ.Networks.payoff, Econ.Networks.IsNash, Econ.Networks.leontief, Econ.Networks.bonacich, Econ.Networks.Dominant, Econ.Networks.remove, Econ.Networks.aggregate, Econ.Networks.intercentrality, Econ.Networks.dominant_pair}
\leanok
Player $i$ earns $\alpha x_i-x_i^2/2+\phi\sum_jg_{ij}x_ix_j$. $M=(I-\phi G)^{-1}$, Bonacich centrality $b=M\mathbf 1$.
\texttt{Dominant}: $g_{ij}\ge0$, $\phi\ge0$ and $\phi\sum_jg_{ij}<1$ for every $i$.
\end{definition}

\begin{lemma}[Properties of $M$ under dominance]
\label{lem:net-dominant}
\lean{Econ.Networks.Dominant.le_of_fixed, Econ.Networks.mulVec_one_sub_apply, Econ.Networks.Dominant.isUnit, Econ.Networks.Dominant.mul_leontief, Econ.Networks.Dominant.leontief_col, Econ.Networks.Dominant.inv_nonneg, Econ.Networks.Dominant.one_le_inv_diag, Econ.Networks.Dominant.one_le_bonacich}
\leanok
\uses{def:net-game}
Solutions of $x=c+\phi Gx$ with $c\ge0$ satisfy $x\ge c$. Hence $I-\phi G$ is invertible, $M\ge0$, $M_{ii}\ge1$, $b\ge1$.
\end{lemma}
\begin{proof}
\leanok
At a minimal coordinate $x_{i}<0$ we would get $x_{i}\ge\phi(\sum_jg_{ij})x_{i}>x_{i}$. So $x\ge0$ and then $x\ge c$.
Injectivity: $(I-\phi G)v=0$ gives $v\ge0$ and $-v\ge0$. Column $k$ of $M$ solves $x=e_k+\phi Gx$.
\end{proof}

\begin{theorem}[BCZ Theorem 1]
\label{thm:net-nash}
\lean{Econ.Networks.payoff_update, Econ.Networks.isNash_iff, Econ.Networks.isNash_iff_eq_bonacich, Econ.Networks.nash_pos}
\leanok
\uses{lem:net-dominant}
With no self-loops, best responses are $x_i=\alpha+\phi\sum_jg_{ij}x_j$, and under dominance the unique Nash
equilibrium is $x=\alpha b$, positive when $\alpha>0$.
\end{theorem}
\begin{proof}
\leanok
Payoff after deviating to $y$ is $-(y-r_i)^2/2+r_i^2/2$ with $r_i=\alpha+\phi\sum_jg_{ij}x_j$ (the $g_{ii}$ term vanishes).
The first-order system is $(I-\phi G)x=\alpha\mathbf 1$.
\end{proof}

\begin{theorem}[BCZ Theorem 3: key player]
\label{thm:net-key}
\lean{Econ.Networks.sum_ne, Econ.Networks.inv_remove, Econ.Networks.aggregate_remove, Econ.Networks.intercentrality_symm, Econ.Networks.keyPlayer, Econ.Networks.Dominant.remove, Econ.Networks.isNash_remove, Econ.Networks.leontief_pair, Econ.Networks.pair_numbers}
\leanok
\uses{thm:net-nash}
Removing $i$ updates $M$ by $M^{-i}_{jk}=M_{jk}-M_{ji}M_{ik}/M_{ii}$, so aggregate centrality falls by the
intercentrality $c_i=b_i\sum_jM_{ji}/M_{ii}$ ($=b_i^2/M_{ii}$ for symmetric $G$). The player whose removal minimizes
equilibrium aggregate effort maximizes $c_i$ (the reduced game is again dominant, with equilibrium $\alpha b(G_{-i})$).
Two linked players with $\phi=1/2$: $M=\begin{pmatrix}4/3&2/3\\2/3&4/3\end{pmatrix}$, $b=(2,2)$, aggregate $4$,
$c_i=3$, aggregate $1$ after removal.
\end{theorem}
\begin{proof}
\leanok
Check $(I-\phi G)_{-i}$ times the proposed matrix entrywise: the terms with index $i$ cancel because
$(I-\phi G)M=I$ and $j\ne i$. Sum the update over $j,k\ne i$; the $i$-th row and column terms vanish identically.
For symmetric $G$, $M$ is symmetric.
\end{proof}
